% Part: first-order-logic
% Chapter: tableaux
% Section: derivations

\documentclass[../../../include/open-logic-section]{subfiles}

\begin{document}

\iftag{FOL}
      {\olfileid{fol}{tab}{der}}
      {\olfileid{pl}{tab}{der}}

\olsection{\usetoken{P}{tableau}}

\begin{explain}
We hebben aangegeven wat een aanname is en de afleidingsregels gegeven.
Uit deze gegevens worden !!{tableau}s inductief voortgebracht: elk
!!{tableau} bestaat uit één tak met een of meer aannamen, of ontstaat
uit !!a{tableau} door op een tak een van de afleidingsregels toe te
passen.
\end{explain}

\begin{defn}[!!^{tableau}]
!!^a{tableau} voor de aannamen $\sFmla{S_1}{!A_1}$, \dots,
$\sFmla{S_n}{!A_n}$ (waarbij elke $S_i$ gelijk is aan $\True$
of~$\False$) is een eindige boom van !!{signed formula}s die aan de
volgende voorwaarden voldoet:
\begin{enumerate}
\item De $n$ bovenste !!{signed formula}s van de boom zijn
  $\sFmla{S_i}{!A_i}$, onder elkaar.
\item Iedere !!{signed formula} in de boom die niet een van de aannamen
  is, ontstaat door een correcte toepassing van een afleidingsregel op
  !!a{signed formula} in de tak erboven.
\end{enumerate}
Een tak van !!a{tableau} is precies dan \emph{gesloten} als hij zowel
$\sFmla{\True}{!A}$ als~$\sFmla{\False}{!A}$ bevat, en anders
\emph{open}. !!^a{tableau} waarin elke tak gesloten is, is een
\emph{gesloten !!{tableau}} (voor zijn verzameling aannamen). Als
!!a{tableau} niet gesloten is, dat wil zeggen als het ten minste één
open tak bevat, is het \emph{open}.
\end{defn}

\begin{ex}
  Elke verzameling aannamen vormt op zichzelf !!a{tableau}, maar is in
  het algemeen niet gesloten. (Zij is vanzelfsprekend alleen gesloten
  als de aannamen al een paar !!{signed formula}s bevatten,
  $\sFmla{\True}{!A}$ en~$\sFmla{\False}{!A}$.)

  Uit !!a{tableau} (open of gesloten) kunnen we een nieuw, groter
  tableau verkrijgen door een van de afleidingsregels toe te passen op
  !!a{signed formula}~$!A$ daarin. De regel voegt een of meer
  !!{signed formula}s toe aan het einde van iedere tak die het voorkomen
  van~$!A$ bevat waarop we de regel toepassen.

  Beschouw bijvoorbeeld de aanname $\sFmla{\True}{!A \land \lnot
    !A}$. Het volgende (open) !!{tableau} bestaat alleen uit die
  aanname:
  \begin{center}
    \begin{tableau}{}
      [\sFmla{\True}{\formula{A} \land \lnot \formula{A}}, just=\TAss]
    \end{tableau}{}
  \end{center}
  We verkrijgen hieruit een nieuw !!{tableau} door de regel
  $\TRule{\True}{\land}$ op de aanname toe te passen. Volgens die regel
  mogen we twee nieuwe regels aan het !!{tableau} toevoegen,
  $\sFmla{\True}{!A}$ en $\sFmla{\True}{\lnot !A}$:
  \begin{center}
    \begin{tableau}{}
      [\sFmla{\True}{\formula{A} \land \lnot \formula{A}}, just=\TAss
        [\sFmla{\True}{\formula{A}}, just={\TRule{\True}{\land}[1]},
          [\sFmla{\True}{\lnot \formula{A}}, just={\TRule{\True}{\land}[1]}
          ]
        ]
      ]
    \end{tableau}{}
  \end{center}
  Wanneer we !!{tableau}s opschrijven, noteren we rechts welke regels we
  hebben toegepast (zo betekent $\TRule{\True}{\land} 1$ dat de
  !!{signed formula} op die regel het resultaat is van toepassing van
  de regel $\TRule{\True}{\land}$ op de !!{signed formula} op
  regel~$1$). Dit nieuwe !!{tableau} bevat nu extra
  !!{signed formula}s, maar op slechts één daarvan
  ($\sFmla{\True}{\lnot !A}$) kunnen we een regel toepassen (in dit
  geval de regel $\TRule{\True}{\lnot}$). Dit levert het gesloten
  !!{tableau}
  \begin{center}
    \begin{tableau}{}
      [\sFmla{\True}{\formula{A} \land \lnot \formula{A}}, just=\TAss
        [\sFmla{\True}{\formula{A}}, just={\TRule{\True}{\land}[1]}
          [\sFmla{\True}{\lnot \formula{A}}, just={\TRule{\True}{\land}[1]}
            [\sFmla{\False}{\formula{A}}, just={\TRule{\True}{\lnot}[3]}, close]
          ]
        ]
      ]
    \end{tableau}{}
  \end{center}
\end{ex}

\end{document}
